\clearpage
\begin{trapbox}
\textbf{易错一：任意镜面都套用 $-k$；把反射又误当成垂直}
\begin{center}
\begin{tikzpicture}[x=1cm,y=1cm]
\draw[p055 frame] (0,0) rectangle (13.8,7.6);
\node[anchor=west,font=\small\bfseries] at (.4,7.1) {同一条原线 $k=2$：反射与垂直不能混用};
\foreach \xx/\kind in {3.5/1,10.3/2}{
\begin{scope}[shift={(\xx,3.8)},scale=.88]
  \draw[p055 axis] (-2.8,0)--(2.8,0);
  \draw[p055 axis] (0,-2.5)--(0,2.5);
  \draw[p055 original] (-1.15,-2.3)--(1.15,2.3);
  \node[CStatement,above] at (1,2.65) {$k=2$};
  \ifnum\kind=1
    \draw[p055 mirror] (-2.35,-2.35)--(2.35,2.35);
    \draw[p055 image] (-2.6,-1.3)--(2.6,1.3);
    \draw[p055 wrong] (-1.15,2.3)--(1.15,-2.3);
    \node[orange!85!black,right] at (2.5,1.25) {$\frac12$};
    \node[red!75!black,right] at (1.1,-2.2) {$-2$};
    \node at (0,-3.05) {镜面 $y=x$：正确取倒数};
  \else
    \draw[p055 mirror] (0,-2.5)--(0,2.5);
    \draw[p055 image] (-1.15,2.3)--(1.15,-2.3);
    \draw[p055 wrong] (-2.6,1.3)--(2.6,-1.3);
    \node[orange!85!black,right] at (1.1,-2.2) {$-2$};
    \node[red!75!black,right] at (2.5,-1.25) {$-\frac12$};
    \node at (0,-3.05) {镜面 $x=0$：正确取相反数};
  \fi
\end{scope}}
\node at (6.9,.45) {橙色实线是实际像直线；红色虚线是误用的方向。两图坐标尺度相同。};
\end{tikzpicture}

\end{center}
\textbf{错法：} 关于 $y=x$ 反射，把 $2$ 变成 $-2$。

\textbf{为什么错、怎样改：} 图左的镜面交换横纵，实际像斜率是 $1/2$；$-2$ 是把水平、竖直镜面的规则错搬过来。水平镜面与竖直镜面才都只反转一个增量，得到 $-k$。

\medskip
\textbf{另一种错法：} 图右的原线与像直线斜率为 $2,-2$，就判定它们垂直。

\textbf{为什么错、怎样改：} 两个有限且非零斜率对应的直线垂直，条件为乘积 $-1$；这里乘积为 $-4$。原斜率为 $2$ 时，垂线斜率应为 $-1/2$（红色虚线），它与实际像直线方向不同。

\medskip
\textbf{速查：} 先读镜面方向，再决定变号、取倒数或用通式；若题目要垂直，另查垂直条件。
\end{trapbox}

\clearpage
\begin{trapbox}
\textbf{易错二：斜率符合，就认定关于指定轴对称}
\begin{center}
\begin{tikzpicture}[x=1cm,y=1cm]
\draw[p055 frame] (0,0) rectangle (13.8,7.3);
\node[anchor=west,font=\small\bfseries] at (.4,6.8) {方向相同的两条线，位置仍可能不同};
\begin{scope}[shift={(4.5,3.3)},scale=.95]
  \draw[p055 axis] (-2.4,0)--(3.2,0) node[right] {$x$};
  \draw[p055 mirror] (0,-1.8)--(0,2.6) node[above] {$x=0$};
  \draw[p055 original] (-1.7,-1.7)--(2.35,2.35);
  \draw[p055 image] (-2.3,2.3)--(1.7,-1.7);
  \draw[p055 wrong] (-.4,2.4)--(3.1,-1.1);
  \fill (0,0) circle (1.8pt);
  \fill[red!75!black] (1,1) circle (1.8pt);
  \node[red!75!black,anchor=west] at (1.5,1) {$(1,1)$};
  \node[CStatement,anchor=west] at (4,2) {原线：$y=x$};
  \node[orange!85!black,anchor=west] at (4,1) {真正的像：$y=-x$};
  \node[red!75!black,anchor=west] at (4,0) {错认的像：$y=-x+2$};
\end{scope}
\node[align=center] at (6.9,.65) {两条下降直线斜率同为 $-1$，却不是同一条线。\\红线与原线的交点不在镜面上，无法关于该镜面互为反射。};
\end{tikzpicture}

\end{center}
\textbf{错法：} 两线斜率为 $1,-1$，就认定关于 $x=0$ 对称。

\textbf{为什么错、怎样改：} 橙线与红线斜率相同，位置却不同。相交直线互为反射时，交点必须在镜面上；此例红线与原线交于 $(1,1)$，不满足。核对方向后，还需核对轴上公共点；平行线则核对它们的位置和到镜面的距离。

\medskip
\textbf{易错三：分母为零就写“无解”或“斜率等于无穷大”}

\textbf{错法：} $k=0$ 时硬套 $1/k$，或在通式分母为零时否认像直线存在。

\textbf{为什么错、怎样改：} 斜率分母表示横增量，横增量为零说明直线竖直；它不是一个有限实数斜率。例5已画出 $(4,3)\mapsto(0,5)$，像直线明确存在。原线水平经 $45^\circ$ 镜面变竖直；反过来则变水平，斜率为零。
\end{trapbox}

\clearpage
\begin{trapbox}
\textbf{易错四：漏掉角平分线中的“内角”}
\begin{center}
\begin{tikzpicture}[x=1cm,y=1cm]
\draw[p055 frame] (0,0) rectangle (13.8,8.5);
\node[anchor=west,font=\small\bfseries] at (.4,8) {相同的两条延长线，选择不同射线，内角就会改变};
\foreach \xx/\kind in {3.5/1,10.3/2}{
\begin{scope}[shift={(\xx,4.75)},scale=.95]
  \draw[p055 helper] (-1.6,-2.4)--(1.6,2.4) (-1.6,2.4)--(1.6,-2.4);
  \draw[p055 mirror] (0,-2.3)--(0,2.3);
  \draw[p055 original,-Stealth] (0,0)--(-1.5,-2.25);
  \node[CStatement,left] at (-1.5,-2.25) {$A$};
  \ifnum\kind=1
    \fill[teal!70!black,opacity=.10] (0,0)--(-1.5,-2.25)--(1.5,-2.25)--cycle;
    \draw[p055 image,-Stealth] (0,0)--(1.5,-2.25);
    \node[orange!85!black,right] at (1.5,-2.25) {$B$};
    \node[above right] at (0,-1.35) {$I$}; \fill (0,-1.35) circle (1.6pt);
    \node[align=center] at (0,-3) {$A,B$ 在轴两侧\\竖直轴平分内角};
  \else
    \fill[orange!85!black,opacity=.08] (0,0)--(-1.5,-2.25)--(-1.5,2.25)--cycle;
    \draw[p055 image,-Stealth] (0,0)--(-1.5,2.25);
    \node[orange!85!black,left] at (-1.5,2.25) {$B$};
    \draw[gray,thick] (0,0)--(-2.4,0) node[left] {内角轴};
    \node[align=center] at (0,-3) {$A,B$ 在轴同侧\\竖直轴平分外角};
  \fi
  \fill (0,0) circle (1.8pt);
  \node at (.65,.25) {$P$};
\end{scope}}
\node at (6.9,.6) {两图都有 $k_{PA}=\frac32,\ k_{PB}=-\frac32$；只有左图的边射线关于竖直轴互为镜像。};
\end{tikzpicture}

\end{center}
\textbf{错法：} 仅凭 $k_{PA}=-k_{PB}$，就认定过 $P$ 的竖直轴平分 $\angle APB$ 的内角。

\textbf{为什么错：} 右图把 $PB$ 换成了原像射线的反向射线，斜率没有变，角的内部却移到了左侧。竖直轴仍对称两条\emph{延长线}，却不再把这两条\emph{射线}互相反射。

\textbf{怎样改：} 对三角形中两条不水平且不竖直的边射线，在斜率互为相反数的前提下，核对 $A,B$ 是否位于竖直轴两侧。横向增量异号时，两条射线关于竖直轴互为镜像，该轴平分内角；横向增量同号时，该轴平分外角。

\medskip
\textbf{内心题的额外核对：} $P,A,B$ 不共线，$I$ 位于三角形内部。这样 $PI$ 才是题目所需的内角平分线；不能只由一条斜率等式倒推“某点就是内心”。
\end{trapbox}
